%% kapa: exam=1 round=21 year=2010 period=1 form=A subject=economics subject_ko=경제원론 seq=24 number=64
%% source: Q-Net 감정평가사 기출문제 (cst00309) artlSeq=5024181 file=extracted/1008/5024181/2010년 제21회 감정평가사 1차시험 기출문제/1교시.hwp
%% conversion: kordoc 4.24.51 -> tools/gen1.py (auto); review=proofread-sonnet; flags=none
\begin{question}{64}
甲과 乙만으로 구성된 A국에서 두 사람이 각각 하루 10시간 일하며, X재와 Y재만을 생산한다. 甲은 시간당 X재 2단위 또는 Y재 1단위를 생산할 수 있으며, 乙은 시간당 X재 1단위 또는 Y재 2단위를 생산할 수 있다. 다음 설명 중 옳지 않은 것은?
\begin{choices}
  \item A국의 X재 하루 최대 생산량은 30이다.
  \item A국의 Y재 하루 최대 생산량은 30이다.
  \item A국의 생산가능곡선은 기울기가 -1인 직선 형태를 지닌다.
  \item 두 사람 모두 하루에 5시간씩 X재와 Y재를 생산하는 것은 비효율적이다.
  \item 甲은 X재 생산에, 乙은 Y재 생산에 비교우위가 있다.
\end{choices}
\end{question}
\answer{3}
